% =============================================================
%  Module Description — Digital Design and Computer Architecture II (CS221-F27)
%  Shown on https://wosm.academy as the course page for CS221-F27.
%  Compile with: pdflatex cs221-module-description.tex
%  Write "TODO(owner): ..." where information is still missing;
%  the website highlights it.
% =============================================================
\documentclass[11pt,a4paper]{article}

\usepackage[utf8]{inputenc}
\usepackage[T1]{fontenc}
\usepackage{lmodern}
\usepackage[margin=2.2cm]{geometry}
\usepackage{booktabs}
\usepackage{tabularx}
\usepackage{array}
\usepackage{enumitem}
\usepackage{xcolor}
\usepackage{titlesec}
\usepackage{fancyhdr}
\usepackage[hidelinks]{hyperref}

% ---------- Module details (edit these) ----------------------
% \ModuleCode must equal the course's `code` in src/content/courses/.
\newcommand{\ModuleTitle}{Digital Design and Computer Architecture II}
\newcommand{\ModuleCode}{CS221-F27}
\newcommand{\Programme}{Foundation programme}
\newcommand{\Level}{Semester 3}
\newcommand{\Semester}{Fall 2027}
\newcommand{\Credits}{7 ECTS}
\newcommand{\ModuleLeader}{Dr. Ali Khudiyev}
\newcommand{\Email}{ali.khudiyev@wosm.academy}
\newcommand{\OfficeHours}{Monday--Friday, 10:00--17:00 (Baku time, UTC+4)}
\newcommand{\Prerequisites}{CS102-F26 Digital Design and Computer Architecture I}

% ---------- Styling ------------------------------------------
\definecolor{accent}{RGB}{0,70,130}
\titleformat{\section}{\large\bfseries\color{accent}}{\thesection.}{0.5em}{}
\titlespacing*{\section}{0pt}{1.2em}{0.5em}
\setlist{itemsep=2pt, topsep=4pt}
\setlength{\parindent}{0pt}
\setlength{\parskip}{4pt}

\pagestyle{fancy}
\fancyhf{}
\lhead{\small \ModuleCode\ -- \ModuleTitle}
\rhead{\small Module Description}
\cfoot{\small \thepage}
\renewcommand{\headrulewidth}{0.4pt}

% =============================================================
\begin{document}

\begin{center}
  {\LARGE\bfseries\color{accent} \ModuleTitle}\\[4pt]
  {\large Module Description}
\end{center}
\vspace{6pt}

% ---------- Summary table (not shown on the website: the course page
% shows these details from the macros above) --------------------------
\begin{tabularx}{\textwidth}{>{\bfseries}l X >{\bfseries}l X}
\toprule
Module code   & \ModuleCode    & Credits      & \Credits \\
Programme     & \Programme     & Level        & \Level \\
Semester      & \Semester      & Prerequisites & \Prerequisites \\
Module leader & \ModuleLeader  & Email        & \href{mailto:\Email}{\Email} \\
Office hours  & \multicolumn{3}{l}{\OfficeHours} \\
\bottomrule
\end{tabularx}

% ---------- Sections (shown on the website) ------------------
\section{Module Overview}
This module introduces the organisation and design of modern computer systems,
from the instruction set architecture to the processor implementation, the memory
hierarchy and input/output. It continues Digital Design and Computer Architecture I
(CS102-F26), where you built the 8-bit DX8 computer from logic gates, and moves to
the scale of a real processor: the RISC-V instruction set.

The first part looks at the computer from the programmer's side: data
representation, assembly-language programming in RISC-V, the Von Neumann
architecture, other instruction set architectures, system architecture, and memory
management with caches. The second part goes down to the hardware: Boolean algebra,
adders and multipliers, memory elements and finite automata, and the design of
single-cycle, multi-cycle and pipelined processors, together with modern parallel
systems and the I/O system. The last part treats the design process itself: logic
minimisation with Karnaugh maps and binary decision diagrams, satisfiability (SAT)
as a tool for circuit design and verification, and the physical design of chips.

\section{Aims}
\begin{itemize}
  \item To explain how software instructions are executed by hardware.
  \item To develop an understanding of processor design and performance.
  \item To give practical experience of programming a real instruction set
        architecture (RISC-V) in assembly language, including its calling convention.
  \item To explain how the memory hierarchy, caches, parallelism and input/output
        shape the performance of a computer system.
  \item To introduce the methods and tools used to design, optimise, verify and
        physically implement digital circuits.
\end{itemize}

\section{Learning Outcomes}
On successful completion of this module, students will be able to:
\begin{enumerate}[label=LO\arabic*.]
  \item Describe the main components of a computer system and how they interact.
  \item Write and analyse programs at the assembly-language level.
  \item Explain datapath and control design, including pipelining and hazards.
  \item Evaluate the effect of cache and memory hierarchy on performance.
  \item Design and analyse combinational and sequential circuits, including adders,
        multipliers, memory elements and finite automata.
  \item Apply logic minimisation (Karnaugh maps, binary decision diagrams) and SAT
        solving to circuit design, and describe the main steps of physical design.
\end{enumerate}

\section{Indicative Content}
\begin{itemize}
  \item Performance metrics and Amdahl's law
  \item Data representation and computer arithmetic
  \item Instruction set architecture (RISC-V, compared with other ISAs such as MIPS,
        ARM and x86)
  \item Assembly-language programming, calling conventions and recursion
  \item Von Neumann architecture and system architecture
  \item Memory hierarchy: caches, virtual memory
  \item Boolean algebra; adders and multipliers; memory elements; finite automata
  \item Single-cycle and multi-cycle datapath; control unit
  \item Pipelining, hazards and branch prediction
  \item I/O, buses and interrupts; implementation at the physical layer
  \item Introduction to parallel and multicore architectures
  \item Logic minimisation: Karnaugh maps and binary decision diagrams
  \item Satisfiability (SAT) in circuit design and verification
  \item Physical design: partitioning, placement, routing and timing
\end{itemize}

\section{Weekly Schedule}
Each week has two lectures, a central exercise with worked examples and a tutorial
in which the solutions of the week's exercise sheet are discussed.

\begin{tabularx}{\textwidth}{c X l}
\toprule
\textbf{Week} & \textbf{Topic} & \textbf{Activity} \\
\midrule
1  & Introduction and Data Representation & Lectures + Exercise 1 \\
2  & Assembly Languages I and II & Lectures + Exercise 2 \\
3  & Assembly Example; Von Neumann Architecture & Lectures + Exercise 3 \\
4  & Other ISAs; System Architecture & Lectures + Exercise 4 \\
5  & Memory Management and Caches I and II & Lectures + Exercise 5 \\
6  & Boolean Algebra; Adders & Lectures + Exercise 6 \\
7  & Multiplication; Memory & Lectures + Midterm exam \\
8  & Processor Architecture I and II & Lectures + Exercise 7 \\
9  & Finite Automata; Multi-cycle Processor & Lectures + Exercise 8 \\
10 & Modern Systems and Parallelism I and II & Lectures + Exercise 9 \\
11 & The I/O System; Physical Layer Implementation & Lectures + Exercise 10 \\
12 & Pipelining I and II & Lectures + Exercise 11 \\
13 & Karnaugh Maps; Binary Decision Diagrams & Lectures + Exercise 12 \\
14 & SAT; Physical Design & Lectures + Exercise 13 \\
\bottomrule
\end{tabularx}

\section{Teaching and Learning Methods}
The module is self-paced. All materials are provided on Moodle: recorded lectures
with slides, central exercises with worked examples, weekly exercise sheets with
tutorial sessions, and lab assignments. Assembly programs are written and run in a
RISC-V simulator, and circuits are built in the Logisim-evolution circuit simulator.
Questions are discussed in the module channel on Discord and in office hours.

\begin{tabularx}{0.6\textwidth}{X r}
\toprule
\textbf{Activity} & \textbf{Hours} \\
\midrule
Lectures            & 41 \\
Labs / Tutorials    & 39 \\
Independent study   & 120 \\
\midrule
\textbf{Total}      & \textbf{200} \\
\bottomrule
\end{tabularx}

\section{Assessment}
The lab assignments are RISC-V assembly programs and circuit designs, submitted
weekly. The midterm exam covers Weeks 1--6; the final exam covers the whole module.

\begin{tabularx}{\textwidth}{X c c c}
\toprule
\textbf{Component} & \textbf{Weight} & \textbf{Due} & \textbf{LOs} \\
\midrule
Lab assignments     & 20\% & Weekly  & LO2, LO3, LO5 \\
Midterm exam        & 30\% & Week 7  & LO1, LO2, LO4, LO5 \\
Final exam          & 50\% & End of module & LO1--LO6 \\
\midrule
\textbf{Total}      & \textbf{100\%} & & \\
\bottomrule
\end{tabularx}

\section{Reading List}
\textbf{Essential}
\begin{itemize}
  \item D.\ A.\ Patterson, J.\ L.\ Hennessy: \textit{Computer Organization and
        Design: The Hardware/Software Interface, RISC-V Edition}, 2nd ed.,
        Morgan Kaufmann, 2020.
\end{itemize}
\textbf{Recommended}
\begin{itemize}
  \item J.\ L.\ Hennessy, D.\ A.\ Patterson: \textit{Computer Architecture:
        A Quantitative Approach}, 6th ed., Morgan Kaufmann, 2017.
  \item W.\ Stallings: \textit{Computer Organization and Architecture}, 11th ed.,
        Pearson, 2019.
  \item S.\ L.\ Harris, D.\ Harris: \textit{Digital Design and Computer
        Architecture: RISC-V Edition}, Morgan Kaufmann, 2021.
  \item A.\ B.\ Kahng, J.\ Lienig, I.\ L.\ Markov, J.\ Hu: \textit{VLSI Physical
        Design: From Graph Partitioning to Timing Closure}, 2nd ed., Springer, 2022.
\end{itemize}

\section{Materials and Licence}
Course materials are published under the
\href{https://creativecommons.org/licenses/by-sa/4.0/}{Creative Commons Attribution-ShareAlike 4.0}
licence unless stated otherwise. See \href{https://wosm.academy/licence/}{wosm.academy/licence} for details.

\section{Policies}
\begin{itemize}
  \item \textbf{Attendance:} Attendance is not mandatory. The module is self-paced, so you can study the materials when it suits you.
  \item \textbf{Late submissions:} Work submitted after the deadline is penalised.
  \item \textbf{Academic integrity:} All submitted work must be your own.
        See the \href{https://wosm.academy/code-of-conduct/}{Code of Conduct}.
\end{itemize}

\end{document}
